\documentclass[11pt]{article}
\usepackage[a4paper,margin=22mm]{geometry}
\usepackage{fontspec}
\setmainfont{DejaVu Serif}
\setsansfont{DejaVu Sans}
\usepackage{microtype}
\usepackage{xcolor}
\usepackage{hyperref}
\usepackage{enumitem}
\usepackage{parskip}
\definecolor{Accent}{HTML}{971D32}
\definecolor{Muted}{HTML}{666666}
\hypersetup{colorlinks=true,urlcolor=Accent,linkcolor=Accent}
\setlist{leftmargin=1.6em,itemsep=0.3em,topsep=0.35em}
\setcounter{tocdepth}{2}
\renewcommand{\contentsname}{Contents}
\newcommand{\sectionrule}{\par\vspace{-0.5em}\noindent\textcolor{Accent}{\rule{\linewidth}{0.6pt}}\par}
\begin{document}

\begin{center}
{\sffamily\bfseries\color{Accent} TOEFL WORD OF THE DAY\par}
\vspace{8mm}
{\fontsize{42}{48}\selectfont\bfseries pervasive\par}
\vspace{2mm}
{\Large\sffamily\color{Accent} /pərˈveɪsɪv/\par}
\vspace{2mm}
{\small\sffamily Local audio phrase: pervasive\par}
{\small\sffamily Audio file: audios/pervasive.mp3\par}
\vspace{2mm}
{\large\sffamily\color{Muted} adjective\par}
\vspace{5mm}
{\sffamily widespread throughout a place or system\par}
\vspace{6mm}
{\large Pervasive describes something that spreads through or affects many parts of a place, society, system, or situation.\par}
\vspace{6mm}
\begin{minipage}{0.84\textwidth}
\itshape “Smartphones have had a pervasive influence on modern communication.”
\end{minipage}\par
\vspace{10mm}
\textbf{Date:} August 16th 2026\quad\textbullet\quad \textbf{Level:} B1--mid-B2 TOEFL
\vfill
Created and compiled by \textbf{Amir Kooshky}\par
\vspace{2mm}
\href{https://t.me/KooshkyTOEFL}{Telegram Channel}\quad\textbullet\quad
\href{https://www.instagram.com/kooshkytoefl}{Instagram}\quad\textbullet\quad
\href{https://t.me/KooshkyTOEFL\_pv}{Personal Telegram}
\end{center}
\newpage
\tableofcontents
\newpage

\section{Meanings and usage}
\sectionrule
\subsection{Meaning 1: Widespread throughout a place or system}
\textbf{Part of speech:} adjective

\textbf{IPA:} /pərˈveɪsɪv/

\textbf{Pronunciation phrase:} pervasive

\textbf{Local audio file:} audios/pervasive.mp3

\textbf{Register:} formal; common in academic, social, political, environmental, and professional writing

\textbf{Definition:} spreading through or affecting many parts of a place, group, system, or situation

\textbf{In simpler words:} present almost everywhere

\textbf{Typical pattern:} pervasive + influence, problem, effect, presence, attitude, or sense

\subsubsection{Natural examples}
\begin{enumerate}
  \item Smartphones have had a pervasive influence on modern communication.
  \item The report describes a pervasive sense of distrust among employees.
  \item Air pollution remains a pervasive problem in many large cities.
  \item Researchers found that social inequality was pervasive throughout the region.
  \item A pervasive smell of smoke remained in the building after the fire.
  \item The study examined the pervasive effects of advertising on consumer behavior.
  \item Corruption had become pervasive across several levels of government.
  \item The internet has become a pervasive part of everyday life.
\end{enumerate}

\subsubsection{Nuance and implication}
\textbf{Central nuance:} Pervasive does not simply mean common. It suggests that something has spread deeply or widely enough to be present in many different parts of a system or environment.

\textbf{Usual implication:} The word is often used for problems, influences, attitudes, or effects that are difficult to avoid because they are widely present.

\textbf{Compared with isolated:} An isolated problem exists only in a limited place or situation. A pervasive problem is widespread and affects many different areas or people.

\subsubsection{Grammar notes}
\textbf{How to use it:} Use it before a noun, as in a pervasive influence, or after a linking verb, as in The problem is pervasive.

\textbf{What can follow it:} It commonly describes influence, effect, problem, presence, attitude, belief, culture, smell, atmosphere, or feeling.

\textbf{Common preposition:} When describing where something is widespread, use throughout, across, or in, as in pervasive throughout society.

\textbf{Position in a sentence:} Words such as highly, increasingly, deeply, and extremely can appear before pervasive.

\subsubsection{Useful patterns}
\textbf{Pattern:} a pervasive influence

\textbf{Use:} Use this for an influence that affects many areas of life or behavior.

\textbf{Example:} Mass media has a pervasive influence on public attitudes.

\textbf{Pattern:} a pervasive problem

\textbf{Use:} Use this for a problem that exists widely throughout a society, organization, or system.

\textbf{Example:} Housing insecurity is a pervasive problem in many rapidly growing cities.

\textbf{Pattern:} a pervasive sense of + feeling

\textbf{Use:} Use this when a feeling is shared widely by many people.

\textbf{Example:} There was a pervasive sense of uncertainty after the announcement.

\textbf{Pattern:} be pervasive throughout + place or system

\textbf{Use:} Use this when emphasizing that something exists in many parts of a larger whole.

\textbf{Example:} The practice was pervasive throughout the industry.

\textbf{Pattern:} pervasive effects of + cause

\textbf{Use:} Use this for consequences that spread across many areas.

\textbf{Example:} The researchers studied the pervasive effects of prolonged drought on rural communities.

\subsubsection{Common wrong usage}
\paragraph{Correction 1}
\textbf{Wrong:} The problem is pervasive in only one small department.

\textbf{Correct:} The problem is limited to one small department.

\textbf{Why:} Pervasive suggests that something is widespread, so it is usually inappropriate for something restricted to one small area.

\paragraph{Correction 2}
\textbf{Wrong:} Technology has pervasively influence on education.

\textbf{Correct:} Technology has a pervasive influence on education.

\textbf{Why:} Use the adjective pervasive before the noun influence.

\paragraph{Correction 3}
\textbf{Wrong:} Pollution is a pervasive throughout the city.

\textbf{Correct:} Pollution is pervasive throughout the city.

\textbf{Why:} Pervasive is an adjective, so do not use a before it unless it directly modifies a countable noun.

\paragraph{Correction 4}
\textbf{Wrong:} The smell was very pervasively.

\textbf{Correct:} The smell was very pervasive.

\textbf{Why:} Use the adjective pervasive after was. Pervasively is an adverb.

\section{Word family}
\sectionrule
\subsection{pervade}
\textbf{Part of speech:} transitive verb

\textbf{IPA:} /pərˈveɪd/

\textbf{Definition:} to spread through and be present in every or many parts of something

\textbf{Relationship to the headword:} Pervade is the verb from which pervasive is formed.

\textbf{Grammar pattern:} pervade + place, system, atmosphere, or aspect of life

\textbf{Usage note:} It is formal and often used for smells, feelings, attitudes, influences, and ideas.

\subsubsection{Examples}
\begin{enumerate}
  \item A strong smell of smoke pervaded the building.
  \item A sense of optimism pervaded the organization after the announcement.
  \item Digital technology now pervades many aspects of daily life.
\end{enumerate}

\subsection{pervasiveness}
\textbf{Part of speech:} uncountable noun

\textbf{IPA:} /pərˈveɪsɪvnəs/

\textbf{Definition:} the quality of existing or spreading throughout many parts of a place, system, or situation

\textbf{Relationship to the headword:} This noun names the quality of being pervasive.

\textbf{Grammar pattern:} the pervasiveness of + influence, problem, technology, or practice

\textbf{Usage note:} It is formal and mainly used in academic or analytical writing.

\subsubsection{Examples}
\begin{enumerate}
  \item The study highlights the pervasiveness of digital technology in modern life.
  \item Researchers were surprised by the pervasiveness of the problem across different age groups.
\end{enumerate}

\subsection{pervasively}
\textbf{Part of speech:} adverb

\textbf{IPA:} /pərˈveɪsɪvli/

\textbf{Definition:} in a way that spreads through or affects many parts of something

\textbf{Relationship to the headword:} This adverb describes something that exists or operates in a widespread way.

\textbf{Grammar pattern:} pervasively + present, distributed, integrated, or used

\textbf{Usage note:} It is much less common than pervasive.

\subsubsection{Examples}
\begin{enumerate}
  \item Digital systems are now pervasively integrated into modern workplaces.
  \item The belief was pervasively present across the organization.
\end{enumerate}

\section{Common collocations}
\sectionrule
\subsection{pervasive influence}
\textbf{Applies to:} Widespread throughout a place or system

\textbf{Meaning:} an influence that affects many areas or people

\textbf{Example:} Social media has a pervasive influence on how people receive news.

\subsection{pervasive problem}
\textbf{Applies to:} Widespread throughout a place or system

\textbf{Meaning:} a problem that exists widely across a population or system

\textbf{Example:} Food insecurity remains a pervasive problem in some regions.

\subsection{pervasive effect}
\textbf{Applies to:} Widespread throughout a place or system

\textbf{Meaning:} an effect that reaches many parts of something

\textbf{Pattern:} pervasive effects of + cause

\textbf{Example:} Urbanization has had pervasive effects on local ecosystems.

\subsection{pervasive presence}
\textbf{Applies to:} Widespread throughout a place or system

\textbf{Meaning:} the condition of being present almost everywhere

\textbf{Example:} The pervasive presence of advertising shapes consumer choices.

\subsection{pervasive sense of}
\textbf{Applies to:} Widespread throughout a place or system

\textbf{Meaning:} a feeling shared widely across a group

\textbf{Pattern:} a pervasive sense of + feeling

\textbf{Example:} There was a pervasive sense of anxiety among employees.

\subsection{pervasive culture}
\textbf{Applies to:} Widespread throughout a place or system

\textbf{Meaning:} a set of attitudes or practices that exists throughout an organization or society

\textbf{Pattern:} a pervasive culture of + noun

\textbf{Example:} The investigation revealed a pervasive culture of secrecy within the organization.

\subsection{pervasive attitude}
\textbf{Applies to:} Widespread throughout a place or system

\textbf{Meaning:} an attitude shared by many people

\textbf{Example:} A pervasive attitude of distrust made cooperation difficult.

\subsection{pervasive inequality}
\textbf{Applies to:} Widespread throughout a place or system

\textbf{Meaning:} inequality that exists broadly across a society or system

\textbf{Example:} The report documents pervasive inequality in access to education.

\subsection{pervasive throughout}
\textbf{Applies to:} Widespread throughout a place or system

\textbf{Meaning:} present in many or nearly all parts of something

\textbf{Pattern:} be pervasive throughout + place, organization, or system

\textbf{Example:} The practice became pervasive throughout the industry.

\subsection{increasingly pervasive}
\textbf{Applies to:} Widespread throughout a place or system

\textbf{Meaning:} becoming more widespread over time

\textbf{Example:} Digital surveillance has become increasingly pervasive in public spaces.

\subsection{the pervasiveness of}
\textbf{Applies to:} Widespread throughout a place or system

\textbf{Meaning:} the degree to which something is widespread

\textbf{Pattern:} the pervasiveness of + phenomenon

\textbf{Example:} Researchers examined the pervasiveness of social media use among teenagers.

\textbf{Usage note:} Pervasiveness is the noun form.

\section{Synonyms and antonyms}
\sectionrule
\subsection{Close synonyms}
\subsubsection{widespread}
\textbf{Part of speech:} adjective

\textbf{Applies to:} Widespread throughout a place or system

\textbf{Shared meaning:} Both describe something that exists over a large area or among many people.

\textbf{Difference:} Widespread means existing in many places or among many people. Pervasive often goes further by suggesting that something has spread deeply through many parts of a system or environment.

\textbf{Example:} The policy received widespread public support.

\subsubsection{prevalent}
\textbf{Part of speech:} adjective

\textbf{Applies to:} Widespread throughout a place or system

\textbf{Shared meaning:} Both can describe something commonly found across a population or area.

\textbf{Difference:} Prevalent means common or frequently found in a particular group or place. Pervasive emphasizes how extensively something spreads through or affects the whole environment.

\textbf{Example:} The disease is particularly prevalent among older adults.

\subsubsection{ubiquitous}
\textbf{Part of speech:} adjective

\textbf{Applies to:} Widespread throughout a place or system

\textbf{Shared meaning:} Both describe something that is present very widely.

\textbf{Difference:} Ubiquitous means seeming to be everywhere. Pervasive emphasizes spreading through and influencing many parts of a system or situation.

\textbf{Example:} Smartphones have become ubiquitous in many countries.

\subsubsection{extensive}
\textbf{Part of speech:} adjective

\textbf{Applies to:} Widespread throughout a place or system

\textbf{Shared meaning:} Both can describe something that reaches a large area or scope.

\textbf{Difference:} Extensive means covering a large amount, area, or range. Pervasive specifically suggests that something penetrates or affects many different parts of a whole.

\textbf{Example:} The storm caused extensive damage across the region.

\subsection{Direct or useful antonyms}
\subsubsection{isolated}
\textbf{Part of speech:} adjective

\textbf{Applies to:} Widespread throughout a place or system

\textbf{Opposition:} Isolated describes something limited to a small number of cases or locations, while pervasive describes something spread throughout many parts of a system or area.

\textbf{Example:} Investigators concluded that the failures were isolated incidents rather than evidence of a larger problem.

\subsubsection{localized}
\textbf{Part of speech:} adjective

\textbf{Applies to:} Widespread throughout a place or system

\textbf{Opposition:} Localized means restricted to a particular place or area, while pervasive means widely present across many areas.

\textbf{Example:} The environmental damage remained localized near the factory.

\section{Words learners may confuse}
\sectionrule
\subsection{prevalent}
\textbf{Part of speech:} adjective

\textbf{Why learners confuse them:} Both are formal adjectives used to describe things that are common or widespread.

\textbf{Key difference:} Prevalent means common or frequently occurring. Pervasive suggests something is not only common but spread deeply or widely through many parts of a system or environment.

\textbf{pervasive:} A pervasive culture of mistrust affected every department.

\textbf{prevalent:} Seasonal allergies are prevalent during the spring.

\subsection{pervasive}
\textbf{Part of speech:} adjective

\textbf{Why learners confuse them:} Pervasive and persuasive look and sound somewhat similar but have unrelated meanings.

\textbf{Key difference:} Pervasive describes something widespread throughout a system. Persuasive, which has a similar spelling, means able to convince people.

\textbf{pervasive:} The report identified a pervasive problem with workplace stress.

\textbf{pervasive:} The researcher presented a persuasive argument for changing the policy.

\section{Etymology}
\sectionrule
\textbf{Origin:} Pervasive comes from the Latin verb pervadere, meaning to go through or spread through.

\textbf{Memory link:} Think of something pervasive as spreading through many parts of a place or system.

\section{Practice exercises}
\sectionrule
\subsection{Word family choice}
Choose the best member of the word family for each blank.

\begin{enumerate}
  \item A strong smell of smoke seemed to \_\_\_\_\_ the entire building.
  \item The study examined the \_\_\_\_\_ of digital technology in modern education.
  \item The problem had become \_\_\_\_\_ throughout the organization.
\end{enumerate}

\subsection{Correct or incorrect?}
Decide whether each sentence is correct. Correct the inaccurate ones.

\begin{enumerate}
  \item The organization faced a pervasive culture of distrust.
  \item The problem was pervasive in one isolated office and nowhere else.
  \item Digital technology has had a pervasive influence on modern society.
  \item The report showed that unemployment was very pervasively.
\end{enumerate}

\subsection{Collocation practice}
Complete each sentence with a natural collocation.

\begin{enumerate}
  \item The investigation revealed a pervasive \_\_\_\_\_ of secrecy throughout the organization.
  \item There was a pervasive sense \_\_\_\_\_ uncertainty among the employees.
  \item The practice became pervasive \_\_\_\_\_ the entire industry.
  \item Researchers studied the pervasive \_\_\_\_\_ of advertising on consumer behavior.
\end{enumerate}

\subsection{Complete answer key}
\subsubsection{Word family choice}
\begin{enumerate}
  \item \textbf{pervade.} A strong smell of smoke seemed to pervade the entire building. \emph{Use the verb pervade after to.}
  \item \textbf{pervasiveness.} The study examined the pervasiveness of digital technology in modern education. \emph{Use the noun pervasiveness after the.}
  \item \textbf{pervasive.} The problem had become pervasive throughout the organization. \emph{Use the adjective pervasive after become.}
\end{enumerate}

\subsubsection{Correct or incorrect?}
\begin{enumerate}
  \item \textbf{Correct} \emph{Pervasive correctly describes a culture that exists widely throughout the organization.}
  \item \textbf{Incorrect}. The problem was limited to one isolated office and occurred nowhere else. \emph{Pervasive implies that something is widespread, not restricted to one small location.}
  \item \textbf{Correct} \emph{Pervasive influence is a natural collocation for an influence that reaches many areas of life.}
  \item \textbf{Incorrect}. The report showed that unemployment was very pervasive. \emph{Use the adjective pervasive after was.}
\end{enumerate}

\subsubsection{Collocation practice}
\begin{enumerate}
  \item \textbf{culture.} The investigation revealed a pervasive culture of secrecy throughout the organization. \emph{A pervasive culture is a set of attitudes or practices found widely across a group.}
  \item \textbf{of.} There was a pervasive sense of uncertainty among the employees. \emph{Use a pervasive sense of followed by a feeling or condition.}
  \item \textbf{throughout.} The practice became pervasive throughout the entire industry. \emph{Pervasive throughout emphasizes that something exists in many parts of a larger system.}
  \item \textbf{effects.} Researchers studied the pervasive effects of advertising on consumer behavior. \emph{Pervasive effects are consequences that reach many different areas.}
\end{enumerate}

\vfill
\begin{center}
\small\color{Muted} Created and compiled by \textbf{Amir Kooshky}\\
\href{https://t.me/KooshkyTOEFL}{Telegram Channel}\quad\textbullet\quad
\href{https://www.instagram.com/kooshkytoefl}{Instagram}\quad\textbullet\quad
\href{https://t.me/KooshkyTOEFL\_pv}{Personal Telegram}
\end{center}
\end{document}
